%% ThirdEye — two-page summary of the MSR 2027 submission.
%% Standalone: compiles independently of main.tex, shares refs.bib.
%%   latexmk -pdf summary.tex

\documentclass[sigconf,nonacm]{acmart}

\usepackage{booktabs}
\usepackage{microtype}
\settopmatter{printfolios=true}

\newcommand{\ci}[2]{\,[#1,\,#2]}
\newcommand{\tool}{ThirdEye}

\begin{document}

\title{A False-Alarm Rate Is Not a Property of the Tool}
\subtitle{Two-page summary of a submission to MSR 2027}

\author{Umar Ahmed}
\affiliation{\institution{PES University}\city{Bengaluru}\country{India}}
\email{PES2UG23CS662@pesu.pes.edu}
\author{Anvita Agarwal}
\affiliation{\institution{PES University}\city{Bengaluru}\country{India}}
\email{PES2UG23CS085@pesu.pes.edu}
\author{Tarun Basavaraju}
\affiliation{\institution{PES University}\city{Bengaluru}\country{India}}
\email{PES2UG23CS644@pesu.pes.edu}
\author{V Lakshitha}
\affiliation{\institution{PES University}\city{Bengaluru}\country{India}}
\email{PES2UG23CS665@pesu.pes.edu}

\begin{abstract}
The same 1,154 verdicts from one unchanged tool, scored two defensible ways,
yield recall of 25.5\% or 82.2\%. We build a Solidity benchmark whose safe class
is stratified by \emph{how the safe label was obtained}, and show that four
unreported evaluation choices each move the headline metric further than any
tool-versus-tool difference in the literature. This summary states the claim, the
evidence, and what the work does not yet establish.
\end{abstract}

\keywords{smart contracts, benchmark construction, evaluation methodology}

\maketitle

\section{The problem}

GPTScan~\cite{gptscan} reports a 4.39\% false-alarm rate.
Heimdallr~\cite{heimdallr} measures \emph{the same tool} at 97.50\%.
Twenty-two-fold, with no change to the tool: the difference lies entirely in
which contracts were designated ``not vulnerable.'' Slither's much-cited 10.9\%
\cite{slither} is not even the same quantity --- it is $1-$precision on one
detector class with no true-negative set --- yet LLM tools are informally held
against it.

Berriche et al.~\cite{berriche} establish the general form of this problem:
across eight cybersecurity benchmarks and ten models, one scoring-pipeline choice
moves a score by more than eighty points. Every benchmark in that study is
question-answering, and \textbf{a multiple-choice question cannot produce a false
alarm}. A detection benchmark is a different object: it must \emph{construct} the
class of things that are not bugs, and \emph{decide} what counts as having found
one. Neither construction has a correct answer, both go unreported, and both ---
we find --- move the headline further than the tool does.

\section{What we built}

A balanced benchmark of 2,250 Solidity contracts (1,125 safe : 1,125 vulnerable)
in which the safe half is stratified by label provenance: audited libraries
(OpenZeppelin, Solady), contracts inside a published audit's scope with no
finding filed, and deployed contracts with no public report --- the last being an
inference from silence rather than evidence of safety. To our knowledge this is
the first corpus that lets a false-alarm rate be decomposed by how much the
``safe'' label can be trusted.

Every contract is checkpointed individually and sampling is nested and seeded, so
every table is re-scorable without re-running inference, and a longer run is a
strict superset of a shorter one. The tool under measurement is a council of
single-class LLM specialists pooling confidences as
$1-\prod(1-\mathit{conf}_i)$ at $\tau = 0.925$; 1,154 contracts were scored.

\section{Four findings}

Each experiment changes one evaluation choice with the tool, seed, model digests
and (except RQ4) the machine held fixed. Intervals are Wilson; paired designs use
McNemar.

\begin{table}[h]
\centering
\caption{Each row is one evaluation choice, with the tool unchanged.}
\label{tab:findings}
\small
\begin{tabular}{@{}p{0.40\columnwidth}rr@{}}
\toprule
\textbf{The unreported choice} & \textbf{Effect} & \textbf{$n$} \\
\midrule
How the safe label was obtained & FPR 15.7\% $\rightarrow$ 37.5\% & 603 \\
What counts as a detection & recall 25.5\% $\rightarrow$ 82.2\% & 349 \\
What the model may read & recall 12.5\% $\rightarrow$ 81.2\% & 20 \\
Which machine ran it & agreement 0.735 & 230 \\
\bottomrule
\end{tabular}
\end{table}

\paragraph{RQ1 --- label provenance.}
False alarms are 15.7\%\ci{11.0}{21.9} on audited libraries against
37.5\%\ci{31.6}{43.8} on audit-reviewed contest code. The effect is two-level
rather than a monotonic gradient: the two weaker tiers overlap. The obvious
confound is size, and it runs the wrong way --- the lowest-false-alarm tier has
nearly twice the median contract length (6,649 characters against 3,531), so the
effect is if anything understated.

\paragraph{RQ2 --- the matching rule.}
Re-scoring the same checkpoints, any-slice-positive gives 82.2\% recall and
type-matched gives 25.5\%: a 56.7-point spread with no verdict changed. That
figure decomposes, and we refuse to claim the sum. 158 of 349 labelled contracts
carry only classes for which the council has no specialist. On the 187 where a
specialist exists, the spread is \textbf{31.6 points} --- still larger than any
tool-versus-tool difference we found in the literature. We report 31.6 as the
protocol effect and 56.7 as protocol and coverage jointly.

\paragraph{RQ3 --- what the model may read.}
The runtime truncates over-long prompts \emph{without error}, and abstentions are
excluded from scoring, so a misconfiguration removes rows instead of lowering a
score. 24.0\% of the corpus overflows a 4,096-token window, skewed towards
vulnerable contracts because they are larger. Widening the window moves recall
$12.5\% \rightarrow 81.2\%$ (McNemar $p = 0.016$) and runs 4.01$\times$
\emph{faster}. The safe-class arm is not reported: that run was destroyed by a
backend outage, so RQ3 stands on the vulnerable class only.

\paragraph{RQ4 --- the machine.}
Two hosts with byte-identical model digests agree on 0.735 of verdicts --- one in
four flips. A \texttt{num\_parallel = 1} control excludes batch parallelism
($p = 0.75$). Restoring a 2.7$\times$ larger model \emph{increases} false alarms
$49 \rightarrow 74$ ($p = 0.00008$): capability is not the lever.

\section{What it costs us}

Scored type-matched, our 25.5\% recall sits inside the band of the static
analysers we set out to beat --- Durieux et al.~\cite{durieux} report Mythril at
27\% and Slither at 17\% on the same style of unit. Where Slither answers at all,
it is more precise than we are (32 tp / 2 fp on 46 of 150 contracts, against our
0.712 precision); it is a high-precision, low-coverage instrument and we are a
moderate-precision, full-coverage one. Against GPTScan on 34 shared projects the
intervals overlap and no difference is demonstrated. The project's original
hypothesis --- that a model-diverse council raises precision --- is refuted on the
one early corpus that had a negative class, and adversarial arbitration, the
standard precision mechanism, wins only 4 of 10 held-out splits.

Reported at 82.2\%, the same comparison would have read as a two-fold win over the
entire field combined. That would have been an artifact of protocol, not of
capability, and a reviewer would have found it. The honest version is the
publishable one: it is what makes the comparison with prior work fair for the
first time.

\section{Status and open items}

The submission targets MSR 2027 (abstract 20 October 2026, paper 23 October). Four
items remain. The safe-class arm of RQ3 must be re-run, and until it is, the
paper must state that 29.9\% may be an under-estimate by an unbounded amount ---
the single most attackable sentence in the draft. RQ3's vulnerable arm rests on 20
paired contracts and should be extended. The 70\% label-noise bound is
self-assessed by the tool's own authors and needs an independent second rater over
the same 21 contracts. Two cited works are still to be read in full rather than
in summary.

We also record seven silent failure modes found during this work, six proposed
invariants, and the fact that our own measurement instrument later violated one of
them --- which is the strongest evidence we can offer that these mistakes are not
obvious in retrospect.

\section{Artifact}

Corpus, loaders, per-contract checkpoints for every run, replay tooling and a
chronological decision log are released at
\url{https://github.com/tarunb05/ThirdEye}: 18,078 lines of source over 61
commits, 42 passing regression tests, 3,707 lines of documentation. Every figure
in the paper traces to a committed result file.

\bibliographystyle{ACM-Reference-Format}
\bibliography{refs}

\end{document}
